Los ratios principales se definen operativamente en el Cap.~1 §1.4 y en el Anexo~D. Resumen:

\begin{itemize}
  \item \textbf{\StructuralScore{} (E):} grado de correspondencia estructural OpenAPI $\leftrightarrow$
        BIAN normalizado a $[0,1]$.
  \item \textbf{\SemanticScore{} (S):} similitud de significado entre términos y conceptos.
  \item \textbf{\CoverageScore{} (C):} proporción de elementos BIAN cubiertos por el contrato.
  \item \textbf{\AlignmentScore{}:} combinación ponderada; clasificación Alta ($\geq 0{,}80$),
        Media, Baja, Nula según §1.4.
\end{itemize}

Los indicadores por fase se detallan en la matriz de consistencia (Anexo~C) y en la
\Cref{tab:anexo-d}.

Contra la etiqueta experta, cuando el subconjunto esté cerrado, se reportan precision, recall
y F1 del emparejamiento. Un verdadero positivo es un par del greedy que el lector confirma.
Un falso positivo es un par que el lector rechaza. Un falso negativo es un par que el lector
marcó y el greedy no armó. Esas tres cifras no se inventan aquí: todavía no hay etiquetas.
Sí se reportan, por instancia y por dominio (cuentas, pagos), S, E, C, el score y el nivel.
El control negativo I6 entra como rechazo, no como fallo del parser.
